DONE - JAY export this to PDF and share with:

Jackson, Thomas, Rebecca, and Isaac


A. Formatting, Proofreading

DONE  - By Sat night - Jason 1. Update index terms

2. DONE - Daniel - Rainbow graph? (Figure 5?)
3. Checking Figure /Text order (is the figure listed after/before being mentioned?)

TODO 4. Logan - Appendix formatting- can we make it fit nicely and parallel for the two experiments?

5. DONE - Jason - update acknowledgments

6. DONE Allyn - rename image filenames to be loyd1.png, loyd2.png....
7. DONE Jason - rename author image files first 5 letters of last name



B. Content Updates
Adjusting based on review feedback:
My other remarks are as follows. 

1. Abstract. “Additionally, we discuss the results of our experiments, focusing in particular on their implications with respect to the question of robustness and transferability of evolved hardware circuits.” In my opinion, the above statement is not informative. The authors should explain what their contribution is. Instead, they only inform us of the effects of their contribution. 

JAY: Good feedback! We should state the implications more explicitly in the abstract

AL: Added "Our results show that our framework has the potential evolve circuits that can be used long-term or outside of environments they were evolved in.
JAY: I revised the wording connecting this addition.


2. Abstract. “Results are promising and demonstrate the feasibility of the approach which, in the long run, has the potential to increase the efficiency and lower the power requirements of modern computing and artificial intelligence systems.” Another unclear statement. I do not understand what the “approach” is. The authors propose a framework that can be used for optimization purposes. Thus, “the approach” should be a problem-dedicated optimizer. However, the paper does not seem to propose anything like that. 

JAY: This would be good to spell out more clearly FPGA intrinsic, bitstream level evolution of circuits. 

DG: addressed in the introduction (just above the section mentioned below), but might need revising

AL: Changed to "Results are promising and demonstrate the feasibility of [[intrinsic, bitstream level evolution]] which, in the long run, has the potential to increase the efficiency and lower the power requirements of modern computing and artificial intelligence systems.



3. Introduction “Consistent with the EHD approach, we aim to evolve a final circuit that does not change during it’s operation.” I do not think that a typical researcher from the EA field will be aware of the details of the EHD approach and will understand the consequences of its use. 

JAY: Good feedback- we need to explain this better. Without a space restriction, this should be very doable. 


DG: addressed, but might need revising
DONE -  JAY - reviewed I think it is reasonable



4. A general style remark. In my opinion, the authors should focus on describing their propositions and the consequences of these propositions. They should avoid using an “enthusiastic” style, which is highly subjective and is not a good style for scientific papers. 

For instance: Section II. “Moreover, because Bitstream Evolution is an open-source platform that utilizes relatively inexpensive and readily available hardware components—an iCE40 FPGA and Arduino Nano MCU have a combined cost of less than \$200—it is extremely accessible.” 

JAY: Fair feedback. We can pull back on the enthusiasm and try to keep the language more straight forward..

TODO LM: will run through and tone down the enthusiasm



Is there a graduation scale for hardware accessibility? I agree that based on price, one can tell whether the equipment is cheap or expensive. How do these two terms refer to “extremely accessible equipment”? I do not know. If authors wish to underline that the equipment in question is cheap, they should rather refer to the average salary in a given group of countries. 

Finally, the authors should report when and where the considered equipment was worth 200 dollars.

JAY: These are great points- it would be good for us to make this a more thorough discussion of research accessibility.

TODO - LM try to add details about specific components and costs
JAY -Logan, you might report on the viability of the cheaper sugar unit as a potential even less expensive option if you have had success with it?


DONE - JAY work on language here to make it more universal and listed online suppliers